% id: 27-02
% book: 베이직쎈 미적분1 · 01 함수의 극한
% section: 학교 시험 기출로 실전 감각 UP!
% figure: tikz:fig-27-02
% answer: 
% answer_source: 
% variant_level: 0 (원본 전사)
% 이 파일은 build.mjs 가 items.json 에서 만든다. 고칠 때는 items.json 을 고친다.
\smash{\raisebox{4.5mm}{\dmkichul{베이직쎈 미적분1 27쪽 02번 · 꼭 나와요}}}
\noindent\begin{minipage}[t]{0.58\linewidth}\vspace{0pt}\raggedright $-2\le x\le 2$에서 정의된 함수 $y=f(x)$의 그래프가 오른쪽 그림과 같을 때, 옳은 것만을 \textbf{보기}에서 있는 대로 고른 것은? \end{minipage}\hfill\begin{minipage}[t]{0.4\linewidth}\vspace{0pt}\centering \resizebox{\ifdim\width>\linewidth\linewidth\else\width\fi}{!}{% 27-02(27쪽) — $-2\le x\le 2$ 에서 정의된 $y=f(x)$ 의 그래프(사분원 호 + 선분 + 수평선분 · 빈 점(○)과 채운 점(●)). 스캔 화질이 낮아 TikZ 로 다시 그림.
% 원문에 있는 것만: 눈금 -2, -1, 1, 2 · y 값 1, 2 · 안내 점선 · 라벨 y=f(x), O, x, y.
\begin{tikzpicture}[x=0.66cm, y=0.62cm, line width=0.6pt]
  \draw[->, >=stealth, line width=0.7pt] (-2.6,0) -- (2.85,0) node[below=1pt] {$x$};
  \draw[->, >=stealth, line width=0.7pt] (0,-0.9) -- (0,2.9) node[left=1pt] {$y$};
  \draw[densely dotted, line width=0.4pt] (-2,0) -- (-2,2);
  \draw[densely dotted, line width=0.4pt] (-1,0) -- (-1,2);
  \draw[densely dotted, line width=0.4pt] (1,0) -- (1,2);
  \draw[densely dotted, line width=0.4pt] (2,0) -- (2,1);
  \draw[densely dotted, line width=0.4pt] (-2,2) -- (1,2);
  \draw[densely dotted, line width=0.4pt] (-1,1) -- (1,1);
  \draw (-2,2) arc (90:0:1 and 1);
  \draw (-1,1) -- (1,2);
  \draw (1,1) -- (2,1);
  \fill (-2,2) circle (2.2pt);
  \fill (-1,2) circle (2.2pt);
  \fill (1,2) circle (2.2pt);
  \fill (2,1) circle (2.2pt);
  \draw[fill=white] (-1,1) circle (2.2pt);
  \draw[fill=white] (1,1) circle (2.2pt);
  \node[below=1pt, font=\small] at (-2,0) {$-2$};
  \node[below=1pt, font=\small] at (-1,0) {$-1$};
  \node[below=1pt, font=\small] at (1,0) {$1$};
  \node[below=1pt, font=\small] at (2,0) {$2$};
  \node[below left=-2pt and -4pt, font=\small] at (0,0) {$\pt{O}$};
  \node[right=1pt, font=\small] at (0,2) {$2$};
  \node[below right=-3pt and 1pt, font=\small] at (0,1) {$1$};
  \node[anchor=west, font=\small] at (1.05,2.55) {$y=f(x)$};
\end{tikzpicture}
}\end{minipage}\par
\noindent \bogi{ㄱ. $\displaystyle\lim_{x\to -1} f(x)$의 값이 존재한다.\\ ㄴ. $\displaystyle\lim_{x\to 1} f(x)$의 값이 존재한다.\\ ㄷ. $-1<a<1$인 실수 $a$에 대하여 $\displaystyle\lim_{x\to a} f(x)$의 값이 항상 존재한다.}
\dmchoicesauto{}{ㄱ}{ㄷ}{ㄱ, ㄴ}{ㄱ, ㄷ}{ㄴ, ㄷ}
